% engines: lualatex, pdflatex
% expect-warnings: 0
% expect-errors: 0
\documentclass{article}
\usepackage[margin=1cm]{geometry}
\usepackage{../../latex/pgf-go}
\newcommand\case[1]{\par\medskip\noindent\textbf{#1}\par\nopagebreak}
\pagestyle{empty}
% Black and white stones, a mark on a stone and one on an empty point, an arrow, a line and a region.
\newcommand\board{%
	\stones[b]{A1}\stones[w]{B1}%
	\marks[b, circle]{A1}\marks[triangle]{A3}%
	\arrow{C1}{E1}\line{C3}{E3}%
	\territory[b, draw = region, line]{E2}}
\begin{document}

% goban/scale changes only the sizes; with scale line width it also multiplies the line widths (the manual
% said that scale did it, and it did not). With layout/fit the widths stay as they are.

\case{scale = 1, 2 and 3: the lines keep their width}
\begin{goban}[5x3, scale = 1]\board\end{goban}
\begin{goban}[5x3, scale = 2]\board\end{goban}
\begin{goban}[5x3, scale = 3]\board\end{goban}

\case{scale line width, scale = 1, 2 and 3: the grid, border, stones, marks, arrow, line and territory get thicker}
\begin{goban}[5x3, scale = 1, scale line width]\board\end{goban}
\begin{goban}[5x3, scale = 2, scale line width]\board\end{goban}
\begin{goban}[5x3, scale = 3, scale line width]\board\end{goban}

\case{scale line width in \textbackslash pgfgoset, layout fit: three boards that fill 4cm columns, with the widths of scale = 1}
\pgfgoset{goban = {scale line width}}
\begin{pgfgolayout}[columns = 3, fit, width = 13cm]
	\begin{goban}[5x3]\board\end{goban}
	\begin{goban}[5x3, scale = 0.5]\board\end{goban}
	\begin{goban}[5x3, scale = 4]\board\end{goban}
\end{pgfgolayout}
\end{document}
